\chapter{Les écarts aux protocoles}
\label{chap:09-ecarts}

Un protocole écrit d'avance ne protège un résultat que si chaque changement ultérieur se voit. Ce
chapitre décrit la règle que nous avons suivie, ce qu'un lecteur peut en vérifier et ce qui repose
sur notre parole. Il donne ensuite la liste complète des écarts des trois protocoles et commente
ceux qui pèsent le plus sur les résultats.

\section{Le principe}
\label{sec:ecarts-principe}

Les trois protocoles datent du 8 novembre 2026 : celui du sondage (\file{study/PROTOCOL.md}), écrit
avant toute lecture d'article, celui des simulations (\file{sim/SIM\_PROTOCOL.md}) et celui du volet
GRPO (\file{sim/llm/LLM\_PROTOCOL.md}), écrits avant tout entraînement. Chacun fixe la question, la
population ou les essais, les mesures et les analyses. Tout changement ultérieur s'ajoute en bas du
fichier, dans une liste numérotée d'\enquote{écarts}, jamais en réécrivant le corps du protocole.

Chaque écart donne trois informations : sa date, sa raison et ce que nous avions déjà vu au moment
de le décider. La dernière est la plus importante. Une règle fixée avant de voir un résultat ne peut
pas avoir été choisie pour lui. Une règle fixée après a pu l'être. Les écarts disent donc
\enquote{AVANT} ou \enquote{APRÈS}, en capitales, avec l'objet précis de la comparaison : avant tout
entraînement, après les résultats de PPO et de SAC mais avant tout résultat de CrossQ, après
l'affichage des valeurs au pas 100 mais avant tout calcul de gain. Trois règles complètent le
dispositif. Aucun essai n'est retiré. Un essai qui plante est relancé une fois avec la même graine aléatoire,
puis l'incident est noté. Un résultat se publie quel qu'il soit, une fois l'analyse annoncée.

Les trois protocoles comptent 66 écarts : 29 pour le sondage, 20 pour les simulations, 17 pour le
volet GRPO. Le \cref{tab:ecarts} les résume tous, la \cref{fig:frise-ecarts} les place dans le temps,
du 8 au 21 novembre 2026.

\section{Ce qui repose sur notre parole}
\label{sec:parole}

Un lecteur peut recalculer chaque nombre à partir des résultats commités
(\cref{chap:08-logiciel}) et vérifier chaque décision du sondage contre la citation et la page sur
lesquelles elle s'appuie. Il ne peut pas vérifier seul l'ordre dans lequel les choses ont été
écrites.

Le dépôt est devenu public le 15 novembre 2026. Son premier commit, du 9 novembre à 6 h 07, contient
déjà des résultats à côté des protocoles. Pour tout ce qui précède le 15 novembre, l'ordre
d'écriture repose sur notre parole et sur les dates des commits, que leur auteur aurait pu forger.
Une relecture critique du dépôt public a de plus trouvé plusieurs déclarations
\enquote{avant}/\enquote{après} fausses dans leurs détails. Les corrections sont datées et
publiées : écart 14 des simulations, écart 15 du volet GRPO. Elles montrent que la règle a été
appliquée avec des erreurs : le lecteur en est averti avant d'accorder sa confiance aux autres
déclarations.

Le dépôt fournit trois appuis pour cette confiance. Les journaux d'évaluation de chaque essai sont
commités, y compris ceux des essais arrêtés quand ils n'ont pas été écrasés. Le registre des
tentatives liste chaque lancement d'entraînement. La provenance des lecteurs du sondage (consigne,
modèle, dates, articles écrits) est reconstruite à partir des transcriptions et publiée dans
\file{study/readers/} (écart 28 du sondage).

\section{Les registres des tentatives}
\label{sec:registres}

Le protocole des simulations dit qu'un essai qui plante est relancé une fois et que l'incident est
noté. Jusqu'au 15 novembre, aucun registre ne gardait la trace de chaque lancement. Le registre
reconstruit, \file{sim/logs/attempts\_reconstructed.csv}, s'appuie sur les dates de création des
journaux, les sorties des lanceurs, les versions commitées et nos journaux de session. Il donne une
ligne par tentative : algorithme, tâche, graine aléatoire, numéro de tentative, début, fin, issue, score vu
avant la relance et preuve. L'écart 14 (e) des simulations en tire le bilan : 292 lancements pour
les 250 entraînements de la première phase, 42 tentatives arrêtées (2 seulement déclarées), aucune plantée.
Chaque arrêt a frappé tous les essais en cours à la fois. Aucun
essai n'a été relancé ni écarté après que son score final était connu. Pour PPO PyTorch et PPO JAX, les
relances retrouvent à l'identique les points d'évaluation des tentatives arrêtées. Pour SAC et
CrossQ, les courbes arrêtées ont été écrasées et ne permettent pas cette vérification.

Depuis le 15 novembre, le lanceur \file{queue.sh} garde un journal par tentative, jamais écrasé. Il
ajoute chaque tentative au registre \file{sim/logs/attempts.jsonl} : début, fin, code de sortie,
empreinte SHA-256 du résultat. Le surveillant y inscrit chaque tâche qu'il arrête. Un arrêt par
signal n'est plus relancé automatiquement. Les essais GRPO arrêtés y figurent aussi.

\section{La liste des écarts}
\label{sec:liste-ecarts}

\gentable{ecarts}

\fig{frise-ecarts}{Les protocoles et leurs écarts datés, du 8 au 21 novembre 2026. Une ligne par
protocole (sondage, simulations, volet GRPO). Chaque marque est un écart, placée à sa date. Sa forme
dit s'il a été écrit avant ou après les résultats qu'il concerne.}

\section{Les écarts du sondage}
\label{sec:ecarts-sondage}

\paragraph{L'harmonisation (écarts 1 à 6, 15 et 19)} Les règles de décision du protocole laissaient
des cas limites ouverts. Les écarts 1 à 6, notés pendant la lecture et avant tout pourcentage,
fixent les plus fréquents : X4 exige un réglage décrit pour les références, X7 une quantité de
calcul, Q5 un lien qui mène à du code. Chaque changement se recode depuis la citation, dans un
fichier d'harmonisation (\file{overrides.json}, \file{overrides\_v2.json}) qui donne la raison de
chaque valeur modifiée. L'accord entre lecteurs se calcule toujours sur les lectures brutes, avant
harmonisation. L'écart 15 corrige une faute de notre part, vue après les pourcentages : les règles
d'harmonisation n'avaient été appliquées qu'aux articles de 2026. Les mêmes règles, sans aucune
règle nouvelle, ont été appliquées aux douze articles de 2018 de la calibration. L'écart 19 applique
l'harmonisation de la grille v2 aux deux années, après un audit des lectures brutes et avant tout
pourcentage : 80 changements, chacun avec sa raison. Pour X3, X4 et X7, aucune règle nouvelle n'a
été ajoutée, afin de lire les deux années avec la même grille. Les cas limites restent codés comme
le lecteur les a lus.

\paragraph{La grille v2 et sa planification (écarts 11, 14, 16 et 17)} La grille v1 ne mesurait
qu'une partie des failles. La grille v2, Y1 à Y12, a été écrite le 9 novembre 2026 avant toute
lecture v2 (écart 11), ses règles de décision le même jour (écart 14), puis précisées après un
pilote de trois articles (écarts 16 et 17), toujours avant le codage principal. Le plan de travail
qui précède ces écarts, \file{study/planification\_v2.md}, a été gardé tel quel et rendu public le
15 novembre 2026. Il annonce, faille par faille, quels critères v2 sont chargés de sortir chaque verdict de
l'état \enquote{à mesurer}.

\paragraph{L'élargissement (écart 12)} L'audit de puissance du premier tirage montre que douze
articles de 2018 et quatorze articles de RL classique de 2026 ne suffisent pas à détecter un écart de
trente points. L'écart 12, écrit avant tout tirage et toute lecture des articles ajoutés, élève chacun des
deux échantillons à 50 articles et fixe la comparaison principale : le RL classique de 2026 contre
notre propre lecture de 2018, avec les mêmes grilles. Les articles de la file de 2026 passent par un
tri sur le titre qui écarte les titres présumés de grands modèles de langue. Le
\cref{chap:10-limites} en discute l'effet.

\paragraph{La règle de verdict (écart 20)} C'est l'écart du sondage qui pèse le plus sur les
conclusions. Il a été décidé après avoir vu les premiers résultats de la lecture élargie. Jusque-là,
un critère avait changé quand l'écart tenait contre les chiffres de Pineau. La comparaison
secondaire montre que notre grille lit 2018 autrement qu'elle : 55~\% contre 90~\% pour les
hyperparamètres de la méthode (Q3), 33~\% contre 10~\% pour l'évaluation tenue à part (Q7). Comparer
2026 à ses chiffres mêlerait le changement des pratiques et celui de la lecture. Depuis cet écart,
un verdict repose seulement sur la comparaison principale de l'écart 12, fixée avant toute lecture
des articles ajoutés. Les comparaisons avec Pineau restent montrées comme contexte. Le changement
de règle est déclaré comme décidé après les résultats. La comparaison qu'il retient avait été fixée
avant eux.

\paragraph{La relecture humaine abandonnée (écarts 13 et 23)} L'écart 13 prévoyait que Warith
Harchaoui relise seul dix articles de 2026, à l'aveugle, sur sept critères. L'écart 23 l'abandonne
le 10 novembre 2026, avant toute relecture. Aucun humain n'a donc relu les articles. Chaque décision
reste vérifiable par sa citation. Le dossier \file{study/relecture\_humaine/} est gardé comme trace
du plan.

\paragraph{Les analyses de réponse à une relecture (écarts 21, 22, 24 à 29)} Après les résultats,
plusieurs analyses ont été ajoutées pour répondre à des objections, chacune écrite avant son propre
calcul et publiée quelle qu'elle soit : les coupes de robustesse et la mise à l'épreuve de la thèse
du coût (écart 21), les listes de contrôle déclarées par les auteurs (écart 22), le filtre de titre
(écart 24), la trajectoire annuelle et sa seconde lecture (écarts 25 et 26), la recherche mécanique
contre les décisions \enquote{non} (écart 27), la provenance des lecteurs (écart 28) et le retrait des
treize articles lus en partie par un autre modèle (écart 29). L'écart 21 a retiré une lecture du
document : la thèse \enquote{Pineau écoutée là où cela coûtait peu} ne résiste pas au test de
corrélation qu'il annonçait (\cref{chap:04-sondage-resultats}).

\section{Les écarts des simulations}
\label{sec:ecarts-sim}

\paragraph{Le budget (écart 1)} SAC et CrossQ s'entraînent 200\,000 pas au lieu de 300\,000, décidé
avant le premier entraînement pour tenir le calcul. PPO garde son million de pas. Cet écart crée la
comparaison à budgets inégaux que l'écart 12 (iii) corrige ensuite par une lecture à budget égal.

\paragraph{B3, écrit en cours de route (écart 8)} L'expérience B3, qui demande si la mesure (score
final, meilleur score, moyenne de la courbe) change le classement, a été écrite le 10 novembre 2026
après les résultats de PPO et de SAC. Elle n'entraîne rien : elle relit les courbes déjà enregistrées.
L'écart 8 la datait \enquote{AVANT tout résultat de CrossQ}. L'écart 14 (b) corrige ce compte, sept
essais CrossQ étant finis à l'écriture. Nos journaux de session, non publics, montrent qu'aucun score
CrossQ n'avait été lu. La formulation exacte devient \enquote{avant d'avoir regardé le moindre
résultat de CrossQ}. L'écart 7, qui fixe les détails de E2, E4 et E5, reçoit la même correction.

\paragraph{Le classement descriptif sous physique modifiée (écart 11)} Le renversement le plus cité
de l'article de blog, PPO écrit en JAX premier sur \task{Hopper} dans la physique d'entraînement et dernier
dans la physique modifiée, vient d'un complément décidé après avoir vu les résultats de E7, sans
test (\file{analysis/e7\_rank.py}). Le document le présente comme tel. La réévaluation par degrés de
l'écart 12 (ii), écrite avant tout nouveau calcul, a ensuite mis ce renversement à
l'épreuve : PPO écrit en JAX n'est plus premier sur \task{Hopper} dans 10 physiques modifiées sur 12.

\paragraph{La consolidation (écart 12)} Écrit le 15 novembre 2026, avant tout nouvel entraînement et
toute nouvelle évaluation, cet écart ajoute trois analyses avec le même code et les mêmes réglages :
vingt graines aléatoires par couple (160 entraînements), la réévaluation de chaque politique sous
des masses et des frottements modifiés par degrés, enfin la lecture de PPO à 200\,000 pas, au budget de
SAC et de CrossQ. Les résultats à dix graines aléatoires restent les résultats principaux. Le
\cref{chap:06-experiences-resultats} donne les deux côte à côte et dit chaque divergence.

\paragraph{La graine aléatoire du bootstrap (écart 13)} En préparant le dépôt public, six scripts ne
redonnaient pas leurs intervalles à l'identique. rliable ignore son argument \texttt{random\_state}
et tire dans le générateur global de NumPy. Chaque script fixe désormais ce générateur au départ ;
deux passages donnent des fichiers identiques à l'octet près. Les estimations ponctuelles et les
verdicts ne changent pas. Les bornes bougent de l'ordre du bruit de Monte-Carlo, au plus 16 points
sur une moyenne interquartile de 2\,423 pour E6.

\paragraph{Sept corrections après une relecture critique (écart 14)} Une relecture hostile du dépôt
public a établi sept points, tous sans effet sur les verdicts.
(a) Le commit qui disait le script de dose \enquote{écrit avant les résultats de dose} se trompait :
30 réévaluations sur 160 existaient déjà. Le script applique une règle écrite avant tout calcul, sans
paramètre libre.
(b) Les comptes d'essais CrossQ finis des écarts 7 et 8 étaient faux.
(c) Le protocole annonçait Stable-Baselines3 2.6.0. Tous les entraînements ont tourné avec la
version 2.9.2, installée avant le premier entraînement.
(d) PPO et TRPO collectent par paquets de 2\,048 pas, si bien que chaque bloc de 50\,000 pas en fait
51\,200 : les modèles de PPO et de TRPO comptent 1\,024\,000 pas, contre 200\,000 pour SAC et CrossQ.
(e) Les relances non déclarées, décrites à la \cref{sec:registres}.
(f) Six scripts lisaient tous les fichiers présents. Avec les graines aléatoires 10 à 19, E6 serait passé de
dix à vingt graines aléatoires sans le dire. Chaque script nomme désormais ses graines aléatoires.
(g) Le registre \file{attempts.jsonl} et les journaux par tentative, à partir du 15 novembre.

\paragraph{L'organisation du calcul (écarts 9, 10, 15 à 17 et 19)} Ces écarts notent des arrêts du
surveillant, des relances, le passage à huit cœurs puis à deux, une suspension liée à la batterie
du Mac, levée le même matin. Aucun ne change un résultat : chaque essai tourne sur un seul
fil et dépend de sa seule graine aléatoire.

\paragraph{Le calcul des ajouts (écart 20)} Écrit le 21 novembre 2026, après la fin de tous les
entraînements et avant tout calcul sur les graines aléatoires 10 à 19, cet écart corrige un défaut de
\file{analysis/e5.py}, qui prenait pour ordre \enquote{correct} la moyenne des dix premiers essais
même avec vingt graines aléatoires. Il liste ensuite les calculs, tous dans \texttt{make analysis}. Deux
d'entre eux, le classement descriptif et la correction de Bonferroni à vingt graines aléatoires, ne figuraient
pas dans l'écart 12 : ils appliquent des règles déjà écrites et sont publiés quels qu'ils soient.

\section{Les écarts du volet GRPO}
\label{sec:ecarts-llm}

\paragraph{Le plan réduit puis rétabli (écarts 6 et 14)} L'écart 6, décidé avant tout essai retenu,
réduit le plan à 100 pas et trois graines aléatoires, fixe le taux d'apprentissage à $3 \times 10^{-6}$ sans
pilote complet et découpe la mise à jour en sous-lots. L'écart 14, écrit le 15 novembre 2026 avant
tout nouvel essai, ajoute les graines aléatoires 3 et 4. Ce rétablissement est la décision qui a le plus changé
la lecture du volet : le gain de Qwen avec la vraie récompense, net à trois graines aléatoires, ne l'est plus à
cinq (\cref{chap:07-modeles-de-langue}).

\paragraph{Les plantages et leur cause (écarts 7 à 11)} Les premiers essais retenus ont planté sur
des probabilités non finies ou manqué de mémoire GPU. Les écarts 7 et 8 posent une règle de relance,
décidée avant de voir un résultat de ces essais. L'écart 9 déplace les essais vers la nuit. L'écart 10
trouve la cause, un calcul GPU interrompu par macOS qui revient avec des valeurs corrompues. Il ajoute
les garde-fous de \file{grpo.py}. Il remet aussi à zéro le compte des échecs, puisque les plantages
passés venaient de la machine : ce changement de règle est décidé après avoir vu ces plantages, alors
qu'aucun résultat de ces essais n'existait. L'écart 11 fait passer les essais Qwen restants sur CPU.
C'est l'origine du mélange entre matériel et graine aléatoire discuté à la
\cref{sec:llm-materiel}.

\paragraph{Corrections après relecture (écart 15)} Écrit le 15 novembre 2026, après les résultats,
il corrige la répartition des neuf échecs du 9 novembre, la date d'un essai, fixe la révision exacte
des deux modèles dans \file{grpo.py}, documente les arrêts du 15 novembre et promet la comparaison
des seuls essais CPU, puisque l'appareil change entre les graines aléatoires de Qwen.

\paragraph{L'essai qui échoue deux fois (écart 16)} Écrit le 17 novembre 2026, après l'affichage,
lors d'un point d'avancement, des valeurs au pas 100 des sept autres nouveaux essais, mais avant tout
calcul de gain ou d'intervalle à cinq graines aléatoires. L'essai Llama à récompense tirée au sort et graine aléatoire 4
échoue deux fois sur le garde-fou du gradient. Selon l'écart 14, ce second échec est un résultat. La
configuration se résume sur ses quatre essais finis, avec une analyse de sensibilité au gain manquant.
L'écart prévoit aussi, hors protocole, un essai unique sur CPU, dont le résultat reste à part.

\paragraph{L'ordre de l'essai sur CPU (écart 17)} Avec le plafond de deux cœurs de l'écart 19 des
simulations, l'essai Llama sur CPU part après tous les autres calculs. Cet ordre ne change aucun résultat.

\section{Ce que ces écarts disent du travail}
\label{sec:ecarts-bilan}

La plupart des écarts tiennent à trois causes : le calcul disponible sur un seul portable, nos
propres erreurs, des objections de relecture. Les écarts décidés après les résultats sont dits tels.
Les trois qui pèsent le plus sur une conclusion sont la règle de verdict du sondage (écart 20), le
classement descriptif sous physique modifiée (écart 11 des simulations) et la remise à zéro du compte
des échecs GRPO (écart 10). Pour le premier, la comparaison retenue avait été fixée avant les
résultats. Le deuxième a été mis à l'épreuve par une analyse écrite d'avance. Le troisième ne touche
aucun nombre publié, puisqu'aucun essai planté n'avait de résultat. Un lecteur qui refuse ces trois
décisions peut refaire chaque calcul sans elles à partir des fichiers commités.
